% RAG 工作流管道图
% 文档 → 分块 → Embedding → 向量库 → 检索 → LLM → 回答
% 上半：离线索引管道  /  下半：在线查询管道
\documentclass[border=5pt,tikz]{standalone}
\usepackage{tikz}
\usetikzlibrary{shapes.geometric, arrows.meta, positioning, fit, backgrounds, calc, decorations.pathreplacing}
\usepackage{xeCJK}
\setCJKmainfont{STSong}

\begin{document}
\begin{tikzpicture}[
    % 索引节点
    node distance=1.5cm,
    idxnode/.style={rectangle, draw=orange!70!black, fill=orange!6, rounded corners=4pt,
                    minimum width=2.4cm, minimum height=1.0cm, font=\footnotesize\bfseries, align=center},
    % 查询节点
    qnode/.style={rectangle, draw=blue!70!black, fill=blue!6, rounded corners=4pt,
                  minimum width=2.4cm, minimum height=1.0cm, font=\footnotesize\bfseries, align=center},
    % 数据节点
    datanode/.style={rectangle, draw=green!60!black, fill=green!6, rounded corners=4pt,
                     minimum width=2.2cm, minimum height=0.9cm, font=\footnotesize, align=center},
    % 箭头
    idxarr/.style={-{Stealth[scale=0.8]}, line width=1.2pt, draw=orange!60!black},
    qarr/.style={-{Stealth[scale=0.8]}, line width=1.2pt, draw=blue!60!black},
    % 标签
    steplabel/.style={font=\tiny\bfseries, above,text=black!50, circle, draw=black!20, inner sep=2pt},
]

% ==============================
% 上半：离线索引管道 (Indexing)
% ==============================
\node[idxnode] (docs) {内部文档\\\footnotesize PDF / Wiki / Log};
\node[idxnode] (chunk) [right=of docs] {文档分块\\\footnotesize Chunking};
\node[idxnode] (embed) [right=of chunk] {向量化\\\footnotesize Embedding};
\node[idxnode] (store) [right=of embed] {向量数据库\\\footnotesize FAISS / Qdrant};
\node[font=\bfseries\small, text=orange!70!black] (indextitle) [above= 0cm of embed] {离线索引管道 (Indexing)};

\draw[idxarr] (docs) -- (chunk) node[midway, steplabel] {1};
\draw[idxarr] (chunk) -- (embed) node[midway, steplabel] {2};
\draw[idxarr] (embed) -- (store) node[midway, steplabel] {3};

% Chunking 细节标注
\node[font=\tiny, text=black!50, align=center, anchor=north] at (chunk.south) {
    \footnotesize 按模块/章节/段落\\
    \footnotesize 父-子块策略
};

% Embedding 细节
\node[font=\tiny, text=black!50, align=center, anchor=north] at (embed.south) {
    \footnotesize BGE-M3 / text2vec\\
    \footnotesize 768$\sim$1024 维向量
};

% ==============================
% 下半：在线查询管道 (Query)
% ==============================
\node[font=\bfseries\small, text=blue!70!black, yshift=0.5cm] (qtitle) [below=of embed] {在线查询管道 (Query)};

\node[qnode] (user) [below=of docs] {用户提问\\\footnotesize Engineering Q};
\node[qnode] (qembed) [right=of user] {问题向量化\\\footnotesize Embedding};
\node[qnode] (search) [right=of qembed] {相似检索\\\footnotesize Top-K};
\node[qnode] (llm) [right=of search] {LLM 生成\\\footnotesize 增强回答};

\draw[qarr] (user) -- (qembed) node[midway, steplabel,] {A};
\draw[qarr] (qembed) -- (search) node[midway, steplabel] {B};
\draw[qarr] (search) -- (llm) node[midway, steplabel] {C};

% ==============================
% 连接：向量库 → 检索
% ==============================
\draw[->, line width=1.5pt, draw=purple!60!black, densely dashed]
    (store.south) -- (search.north)
    node[midway, right, font=\tiny\bfseries, text=purple!60!black] {检索};

% ==============================
% 数据源框
% ==============================
\node[datanode, anchor=south] (data1) [above=of docs] {IP 手册};
\node[datanode, anchor=south] (data2) [above=of chunk] {Bug DB};
\node[datanode, anchor=south] (data3) [above=of embed] {Spec 文档};
\node[datanode, anchor=south] (data4) [above=of store] {Tapeout 经验};

% ==============================
% 背景框
% ==============================
\begin{scope}[on background layer]
        \node[draw=orange!30, fill=orange!3, rounded corners=8pt, inner sep=6pt,
            fit=(docs)(chunk)(embed)(store)] (idxbox) {};
        \node[draw=blue!30, fill=blue!2, rounded corners=8pt, inner sep=6pt,
            fit=(user)(qembed)(search)(llm)] (qbox) {};
\end{scope}

% ==============================
% 底部注释
% ==============================
\node[font=\footnotesize, text=black!60, align=left, ] (foot) [below = 0cm of qbox] {
    \textbf{核心价值：} LLM 基于检索到的真实文档回答，而非"参数记忆"。\\
    在 IC 设计中，一个编造的寄存器地址足以导致芯片失效。
};

\end{tikzpicture}
\end{document}
